  \documentclass[final]{article}
% Landscape poster canvas (A0 landscape)
\usepackage[paperwidth=118.9cm,
paperheight=84.1cm,
margin=1.2cm,
nohead, nofoot, nomarginpar]{geometry}
\setlength{\pdfpagewidth}{118.9cm}
\setlength{\pdfpageheight}{84.1cm}

\usepackage[T1]{fontenc}
\usepackage[utf8]{inputenc}
\usepackage{lmodern}
% Body font: Open Sans (set as the document sans-serif default)
\usepackage{opensans}
\usepackage{helvet} % provide Helvetica family for headings
\renewcommand{\familydefault}{\sfdefault} % make Open Sans the default family
\usepackage{microtype}
\usepackage{amsmath,amssymb}
\usepackage{booktabs}
\usepackage{array}
\usepackage{tabularx}
\usepackage{enumitem}
\usepackage{graphicx}
\usepackage[table]{xcolor}
\usepackage{tikz}
\usepackage{multicol}
\usepackage[most]{tcolorbox}
\usepackage{ragged2e}
\usepackage{hyperref}

\pagestyle{empty}
\setlength{\parindent}{0pt}
\setlength{\parskip}{0pt}           % Reduced global parskip
\setlength{\columnsep}{1.0cm}
\setlength{\tabcolsep}{5pt}
\renewcommand{\arraystretch}{1.14}
\RaggedRight

% Reduce vertical space around displayed equations for tighter poster layout
\setlength{\abovedisplayskip}{6pt}
\setlength{\belowdisplayskip}{6pt}
\setlength{\abovedisplayshortskip}{6pt}
\setlength{\belowdisplayshortskip}{6pt}

\setlist[itemize]{leftmargin=*, topsep=0.05em, itemsep=0.02em, parsep=0pt, partopsep=0pt}
\setlist[enumerate]{leftmargin=*, topsep=0.05em, itemsep=0.02em, parsep=0pt, partopsep=0pt}

% Poster-scaled text sizes tuned for readability at A0
\renewcommand{\normalsize}{\fontsize{31}{35}\selectfont}
\renewcommand{\small}{\fontsize{31}{35}\selectfont}
\renewcommand{\footnotesize}{\fontsize{29}{33}\selectfont}
\renewcommand{\scriptsize}{\fontsize{25}{29}\selectfont}

% --- Minimal IBA-inspired palette ---
\definecolor{ibaMaroon}{RGB}{111,24,31}
\definecolor{ibaMaroonDark}{RGB}{86,16,22}
\definecolor{ibaCream}{RGB}{249,245,237}
\definecolor{ibaPaper}{RGB}{255,253,250}
\definecolor{ibaInk}{RGB}{41,34,32}
\definecolor{ibaGrey}{RGB}{232,228,222}
\definecolor{ibaSoft}{RGB}{244,238,231}
\definecolor{ibaAccent}{RGB}{153,58,58}

% Light red row shade for alternating table rows
\colorlet{rowshade}{ibaAccent!15}

\hypersetup{
	colorlinks=true,
	linkcolor=ibaMaroon,
	urlcolor=ibaMaroon,
	citecolor=ibaMaroon,
	pdfauthor={Uzair Arif and Ahmed Essa},
	pdftitle={Endogenous Institutions for Sustaining Cooperation}
}

\newcommand{\smallhead}[1]{\vspace{0.70em}\noindent{\fontfamily{phv}\selectfont\textbf{\color{ibaMaroon}#1}}\par\vspace{0.08em}}
\newcommand{\mainhead}[1]{\vspace{0.55em}\noindent{\fontfamily{phv}\selectfont\fontsize{36}{40}\selectfont\bfseries\color{ibaMaroon}#1}\par\vspace{0.01em}\noindent\rule{\linewidth}{0.7pt}\vspace{0.08em}}

\newenvironment{posterbody}{\normalsize}{\par}

\begin{document}
	\begin{tcolorbox}[
		enhanced,
		colback=ibaMaroon,
		colframe=ibaMaroon,
		boxrule=0pt,
		arc=0mm,
		width=\textwidth,
		left=4mm, right=4mm, top=7mm, bottom=5mm,
		borderline west={1.4pt}{0pt}{ibaMaroonDark},
		borderline east={1.4pt}{0pt}{ibaMaroonDark}
		]
		\begin{tikzpicture}[remember picture, overlay]
			\node[anchor=west, inner sep=0pt] at (-13cm, -0.25cm) {
				\IfFileExists{Institute_of_Business_Administration,_Karachi_(logo).png}{%
					
\includegraphics[width=3.2cm,height=3.2cm,keepaspectratio]{Institute_of_Business_Administration,_Karachi_(logo).png}%
				}{%
					\fcolorbox{white}{ibaMaroonDark}{%
						\parbox[c][3.5cm][c]{3.5cm}{\centering\color{ibaCream}\bfseries IBA\\LOGO}%
					}
				}
			};
		\end{tikzpicture}
		\centering
		{\fontsize{50}{54}\selectfont\bfseries\color{white} Endogenous Governance of the Loss and Damage Fund: Evidence from LLM-Agent Public Goods Games}\\[0.4em]
		{\fontsize{30}{34}\selectfont\color{white} Uzair Arif (27184) \& Ahmed Essa (26753) \quad Supervisor: Dr. Ghamz-e-Ali Siyal}
	\end{tcolorbox}
	
	\vspace{0.4mm}
	
	\begin{posterbody}
		\begin{multicols}{4}
			\mainhead{1. Background and Motivation}
			\begin{itemize}[leftmargin=*, nosep]
				\item The redistribution of loss and damage finance is debated globally, especially in how support should reach developing countries.
				\item Empirical evidence on how developed countries behave in these settings is still limited.
				\item This study uses LLM-agent public-goods games to predict the outcomes from repeated interaction and institutional design.
			\vspace{2.2cm}
			\end{itemize}
			
			\smallhead{Research Question}
			\begin{itemize}[leftmargin=*, nosep]
				\item Do developed countries sustain cooperation when contributing to a Loss and Damage Fund?
				\item Does the Fund effectively redistribute resources to developing countries during climate crises?
			\end{itemize}
			\vspace{2.2cm}
			
			\mainhead{2. Introduction}
			This study models repeated climate cooperation with endogenous rule changes and budget-constrained contributions. We evaluate whether institutional design can sustain mitigation efforts while providing reliable cross-country support after climate shocks.
			
			In one-shot play, cooperation is difficult to sustain because agents have incentives to under-contribute. In LDF mode, Stage 1 contributions are treasury-constrained (agent-specific wealth budgets), not fixed equal endowments.
			\vspace{2.2cm}
			
			\smallhead{Main Contribution}
			\begin{itemize}[leftmargin=*, nosep]
				\item 26-country repeated climate public-goods game with treasury-constrained contributions and heterogeneous fiscal capacity.
				\item Endogenous governance: periodic constitutional voting updates enforcement and LDF policy parameters.
				\item Joint economic evaluation of mitigation, welfare, inequality, and LDF transfer efficiency.
			\end{itemize}
			\vspace{2.2cm}
			
			\mainhead{3. Literature}
			{\small
				\begin{itemize}[leftmargin=*, nosep]
					\item \textbf{Ostrom (1990)}: Governing the Commons: The Evolution of Institutions for Collective Action.
					\item \textbf{Fehr and G\"achter (2000)}: Cooperation and Punishment in Public Goods.
					\item \textbf{Piedrahita et al. (2025)}: Corrupted by Reasoning: Reasoning Language Models Become Free-Riders in Public Goods Games.
					\item \textbf{Willis et al. (2026)}: Evaluating Collective Behaviour of Hundreds of LLM Agents.
					\item \textbf{Tewolde et al. (2026)}: CoopEval: Benchmarking Cooperation-Sustaining Mechanisms and LLM Agents in Social Dilemmas.
				\end{itemize}
			}
			\vspace{1.5cm}
			
			\mainhead{4. Design and Identification}
			\smallhead{Simulation Design}
			{\small
				\begin{tabularx}{\dimexpr2\columnwidth+\columnsep\relax}{>{\bfseries\raggedright\arraybackslash}p{5.3cm}@{\hspace{1.6cm}}>{\raggedright\arraybackslash}p{\dimexpr\linewidth-5.70cm\relax}}
					Agents          	& 26 Llama 3.1 7B agents (12 developed, 14 developing) \\
					Rounds          & 15 repeated rounds \\
					Stage 1         	& Treasury-based budget cap $0\le c_i^{(t)}\le W_i^{(t)}$, multiplier $\alpha=1.6$ \\
					Initial treasuries & Developed: heterogeneous profile (245, 100, 51, ..., 2); Developing: 3.0 each \\
					Institutions    	& SI (sanction/reward enabled) and SFI (no sanction channel) \\
					Climate shocks  	& Realized in rounds 9 and 13 \\
					LDF mechanism   	& Pool tracked each round; contributions follow replenishment policy; payouts are damage/vulnerability weighted and coverage-capped \\
					Governance      	& Constitutional update rounds at 5, 10, and 15 \\
				\end{tabularx}
			}
			
			\smallhead{Identification Logic}
			\begin{itemize}[leftmargin=*, nosep]
				\item Trace mechanism interactions over 15 rounds under a fixed LDF policy architecture.
				\item Report level, distribution, and transfer metrics jointly to avoid mean-only inference.
				\item Treat results as within-run mechanism evidence, not cross-design causal identification.
				\item Use this poster for core evidence and the report for derivations and robustness checks.
			\end{itemize}
			
			\smallhead{Outcome Variables}
			Contribution path, contribution share of wealth, sanction intensity, welfare, inequality, and LDF transfer efficiency.
			
			\smallhead{Hypotheses}
			\begin{itemize}[leftmargin=*, nosep]
				\item $H_1$: Repeated interaction and enforcement can sustain high mitigation effort.
				\item $H_2$: Shock realization raises contribution effort through salience and loss exposure.
				\item $H_3$: LDF payout rules can leave fund inflows much higher than payouts.
			\vspace{1.1cm}
			\end{itemize}
			
			\mainhead{5. Theory and Mechanism}
			\smallhead{Stage 1 — Contributions / Utility}
			\[
			\pi_i^{(1,t)} = \bigl(W_i^{(t)} - c_i^{(t)}\bigr) + \frac{\alpha\sum_{j=1}^{n} c_j^{(t)}}{n}
			\]
			\vspace{-1mm}
			{\footnotesize LDF-mode budget constraint: $0\le c_i^{(t)}\le W_i^{(t)}$; MPCR $=\alpha/n<1$ for $n>1.6$.}
			
			\vspace{0.6cm}
			\smallhead{Stage 2 — Sanctions / Enforcement}
			\[
			\pi_i^{(2)} = E_2 - \sum_j p_{ij} - \sum_j r_{ij} + \beta \sum_j p_{ji} + \rho \sum_j r_{ji}
			\]
			\vspace{-1mm}
			{\footnotesize $E_2=20$ (SI only), punishment impact $\beta=3$, reward gain $\rho=1$.}
			\par\columnbreak
			\smallhead{Subsidy Redistribution}
			\[
			\sigma_i = \mathbf{1}\{i \in \mathcal{T}\}\; \frac{\phi \sum_k \sum_j p_{kj}}{|\mathcal{T}|}
			\]
			\vspace{-1mm}
			{\footnotesize Fraction $\phi=0.2$ of punishment spending is rebated to top contributors.}
			\vspace{1.0cm}
			
			\smallhead{Round Utility}
			\[
			U_i^{(t)} = \pi_i^{(1,t)} + \pi_i^{(2,t)} + \sigma_i^{(t)}
			\]
			\vspace{-2mm}
			{\footnotesize Cumulative utility captures institutional path dependence.}
			\vspace{1.6cm}
			
			\smallhead{Mechanism Interpretation}
			\begin{itemize}[leftmargin=*, nosep]
				\item $\alpha/n<1$ implies private under-provision despite social gains.
				\item Deterrence works when expected sanction loss exceeds one-shot deviation gain.
				\item Rewards and rebates partly offset enforcement burdens for contributors.
				\item Voting endogenizes constitutional parameters, turning play into institutional adaptation.
			\vspace{1.4cm}
			\end{itemize}
			
			\smallhead{Comparative-Statics Predictions}
			\begin{itemize}[leftmargin=*, nosep]
				\item Larger punishment impact should strengthen deterrence when free-riding pressure rises.
				\item Higher shock salience should increase contribution effort and accelerate coordination.
				\item Weak payout rules should increase LDF accumulation while reducing effective insurance value.
			\vspace{1.5cm}
			\end{itemize}
			
			\mainhead{6. Results and Policy Implications}
			\smallhead{Key Results}
			{\scriptsize
				\setlength{\tabcolsep}{2pt}
				\renewcommand{\arraystretch}{1.1}
				\begin{tabularx}{\linewidth}{>{\bfseries\raggedright\arraybackslash}X >{\raggedleft\arraybackslash}X}
					Metric & Value \\
					\midrule
					Agents / Rounds & 26 / 15 \\
					Total contributions & 481,300 \\
					Mean contribution share of wealth & 8.1\% \\
					Total welfare & 18,095,058.63 \\
					Sanction intensity & 1,686 punish; 1,177 reward \\
					Climate shocks & 2 shocks (R9, R13), 471.31 gross damage \\
					LDF contributions vs payouts & 34,500.37 vs 424.18 \\
					LDF final pool & 34,076.18 \\
					Mean Gini wealth & 0.645 \\
				\end{tabularx}
			}
			
			\smallhead{Key Figures}
			\begin{center}
				\IfFileExists{../fig_ldf_pool.png}{%
					\includegraphics[width=0.98\linewidth,height=12.2cm,keepaspectratio]{../fig_ldf_pool.png}%
				}{% 
					\fbox{\parbox[c][12.2cm][c]{0.95\linewidth}{\centering\textbf{LDF Pool Balance Figure Missing}\\[0.6em]Missing file: ../fig\_ldf\_pool.png}}%
				}
				
				\vspace{0.6cm}
				
				\IfFileExists{../fig_mean_contribution.png}{%
					\includegraphics[width=0.98\linewidth,height=11.2cm,keepaspectratio]{../fig_mean_contribution.png}%
				}{% 
					\fbox{\parbox[c][11.2cm][c]{0.95\linewidth}{\centering Missing image file:\ \texttt{../fig_mean_contribution.png}}}%
				}
			\end{center}
			
			\smallhead{Key Findings}
			\begin{itemize}[leftmargin=*, nosep]
				\item \textbf{Mitigation effort scales sharply}: Mean contribution rises from 5.54 to 7,696.08 million dollars (GDP) between rounds 1 and 15, with the steepest acceleration after shock rounds 9 and 13.
				\item \textbf{Enforcement is economically active}: SI records 1,686 punishment tokens and 1,177 reward tokens over 15 rounds, indicating non-trivial enforcement activity.
				\item \textbf{Redistribution is the binding margin}: 34,500.37 million dollars (GDP) enters the fund, but only 424.18 is disbursed (1.2\%), leaving a terminal pool of 34,076.18 and showing that collection vastly outpaces transfer.
				\item \textbf{Pooling does not automatically equalize outcomes}: Mean Gini remains 0.645 and the developed-to-developing final wealth ratio stays near 5.55x, so the fund accumulates faster than it redistributes.
			\end{itemize}
			
			\smallhead{Interpretation}
			\begin{itemize}[leftmargin=*, nosep]
				\item Agents increase contributions after shocks, showing attention to losses.
				\item There is a lot of enforcement activity alongside high giving.
				\item Most money collected stays in the fund; few payouts happen.
				\item The fund acts like a savings buffer, not an active payer here.
			\end{itemize}
			
			\smallhead{Conclusion}
			Repeated interaction and enforcement can raise cooperation, but the fund only helps if payouts are designed to be used rather than hoarded.
			
			\smallhead{Policy Implications}
			\begin{itemize}[leftmargin=*, nosep]
				\item Link incentives to available treasury so poorer agents are not excluded.
				\item Design payouts and coverage so the fund is actually used.
				\item Add automatic payout rules or triggers to avoid hoarding.
			\end{itemize}
			
			\smallhead{Limitations}
			\begin{itemize}[leftmargin=*, nosep]
				\item \textbf{Institutional assignment is exogenous}: In LDF mode, developed agents are routed to SI and developing agents to SFI in code, so the poster focuses on mechanism performance rather than institution-choice behavior.
				\item \textbf{Single-run evidence}: This poster highlights one focal realization; multi-seed runs would be needed to add uncertainty bands around the magnitudes.
				\item \textbf{Shock process is exogenous}: Climate shocks are parameter-driven stochastic events, so the results describe adaptation under realized shock timing and severity.
			\end{itemize}
			
		\end{multicols}
		
	\end{posterbody}
	
\end{document}